\chapter{Inscrição profissional de pessoa física}%
\label{chap:pf-inscricao}

\section{O que é?}

É o registro necessário para a atuação profissional em Psicologia. O exercício profissional no estado de São Paulo estará autorizado após a geração do número de registro profissional.

\section{Quem deve se inscrever?}

Profissionais graduadas/os em Psicologia, que tenham atuação no estado de São Paulo.

\section{Como solicitar o serviço?}

Por meio do \textit{site} do CRP SP, na guia \href{https://www.crpsp.org/pagina/view/308}{Atendimento > Pessoa Física}.

\section{Que informações ou documentos são necessários?}

\begin{itemize}
\item	 Documento de identificação (RG, CNH etc.);
\item	 Cadastro de Pessoa Física (CPF);
\item	 certidão de quitação eleitoral;
\item	 comprovante ou declaração de endereço;
\item	 diploma de formação em Psicologia (frente e verso) ou certidão de colação de grau;
\item	 certidão de casamento ou de união estável;
\item	 comprovante de quitação com o serviço militar (para homens até 45 anos).
\end{itemize}

\section{Quais são as etapas para a realização desse serviço?}

\begin{enumerate}
\item	A/O psicóloga/o realiza o requerimento de inscrição \textit{on-line}, pelo \href{https://www.crpsp.org/pagina/view/308}{\textit{site} do CRP SP};
\item	o CRP SP faz a análise da documentação;
\item	os boletos para pagamento da taxa de inscrição e da anuidade são emitidos;
\item	após o pagamento das taxas, o registro é deferido e uma declaração para atuação profissional é disponibilizada.
\end{enumerate}

\section{Quanto custa?}

Os valores para 2024 ficaram estabelecidos assim:

\begin{itemize}
\item	pagamento à vista: R\$ 575,60;
\item	pagamento parcelado: 6 x R\$ 95,93.
\end{itemize}

\section{Quanto tempo leva?}

Até 30 dias.

\section{Como ter acesso ao atual \emph{status} do serviço?}

Ao se fazer a solicitação é gerada uma senha de acesso. Com essa senha, é possível acompanhar o \textit{status} da solicitação \textit{on-line}.

\section{Legislação relacionada}

\begin{itemize}
\item	\href{https://atosoficiais.com.br/cfp/resolucao-administrativa-financeira-n-3-2007-institui-a-consolidacao-das-resolucoes-do-conselho-federal-de-psicologia}{Resolução CFP nº~03/2007} (Consolidação das Resoluções do Conselho Federal de Psicologia);
\item	\href{https://transparencia.cfp.org.br/wp-content/uploads/sites/7/2020/05/Manual-de-Procedimentos-Administrativos-e-Financeiro-CFP-RES-20.2018.pdf}{\href{https://atosoficiais.com.br/cfp/resolucao-administrativa-financeira-n-20-2018-altera-o-manual-de-procedimentos-administrativo-e-financeiro-do-sistema-conselhos-de-psicologia-anexo-da-resolucao-cfp-n-202018-a-resolucao-cfp-n-03-2007-a-resolucao-cfp-n-16-2019-e-da-outras-providencias}{Resolução CFP nº~20/2018}}, que determina a revisão e ampliação do \emph{Manual de procedimentos administrativos e financeiros} \href{https://transparencia.cfp.org.br/wp-content/uploads/sites/7/2020/05/Manual-de-Procedimentos-Administrativos-e-Financeiro-CFP-RES-20.2018.pdf}{(versão atualizada)}.
\end{itemize}

\section{Serviços correlatos}

\begin{itemize}
\item	\hyperref[chap:pf-cip]{Entrega da carteira de identidade profissional (CIP)};
\item	\hyperref[chap:pf-inscricao_definitiva]{alteração de inscrição provisória para definitiva};
\item	\hyperref[chap:pf-secundaria]{inscrição secundária};
\item	\hyperref[chap:pf-transferencia]{transferência do registro profissional};
\item	\hyperref[chap:pf-cancelamento]{cancelamento da inscrição};
\item	\hyperref[chap:pf-prorroga_diploma]{prorrogação do prazo de entrega do diploma};
\item	\hyperref[chap:pf-isencao]{solicitação de isenção de anuidade}.
\end{itemize}

\sumario
